\section{A forced-mirror fourth-derivative condition}\label{sec:mirror}

\subsection{Deflation and the exact physical inequality}

At an all-real event \(a>0\) of multiplicity \(m\ge2\), deflate two
copies at \(a\):
\[
 G_a(z)=\frac{H_T(z)}{(z-a)^2}.
\]
The remaining zero multiset still contains at least two copies at
\(-a\). For real \(z\) away from the zeros, logarithmic differentiation
of its product gives
\[
 (G_a'(z))^2-G_a(z)G_a''(z)
  =G_a(z)^2\sum_{\rho:\,G_a(\rho)=0}(z-\rho)^{-2}
  \ge\frac{2G_a(z)^2}{(z+a)^2}.
\]
The inequality extends to \(z=a\) by continuity, including when \(G_a(a)=0\).
Write \(H_j=H_T^{(j)}(a)\). Taylor expansion gives
\[
 G_a(a)=H_2/2,\qquad G_a'(a)=H_3/6,\qquad G_a''(a)=H_4/12.
\]
Multiplication by \(72\) proves the inherited mirror condition
\begin{equation}\label{eq:physicaljet}
 \boxed{2H_3^2-3H_2H_4\ge\frac9{a^2}H_2^2.}
\end{equation}
For \(m\ge4\), both sides vanish; for \(m=3\) the inequality remains
valid with \(H_2=0\). No restriction to exact doubles was made.
The deflation condition and its earlier model compatibility are owned
by \cite{collision}; this section develops its RSCT normalization and
joint calibrated use.

\subsection{Cancellation of the carrier derivatives}

Let \(f_j=\partial_w^jf_a(0)\), with \(f_0=f_1=0\), and evaluate
\(b=A_z/A\) and \(b_z\) at \(z=a\), holding \(T\) fixed.
Differentiating \(H_T=A f_a\) gives
\begin{align}
 H_2/A&=L^2f_2,\\
 H_3/A&=L^3f_3+3bL^2f_2,\\
 H_4/A&=L^4f_4+4bL^3f_3+6(b^2+b_z)L^2f_2.
\end{align}
Substitute in \eqref{eq:physicaljet}. The \(bf_2f_3\) and \(b^2f_2^2\)
terms cancel exactly. After division by \(A(T,a)^2L^6>0\), the condition is
\begin{equation}\label{eq:Jdef}
 \cJ_a(f_a):=2f_3^2-3f_2f_4-\mu_a f_2^2\ge0,
 \qquad \mu_a=\frac{18b_z(a)+9/a^2}{L^2}.
\end{equation}
From \eqref{eq:bcarrier},
\[
 b_z(a)=-\frac7{4a^2}+\frac{T(1-L)}{8a^2},
\]
and hence
\begin{equation}\label{eq:mu}
 \mu_a=\frac{-45/2-9c/4+9c/(4L)}{a^2L^2}<0\qquad(L\ge1).
\end{equation}
The normalization therefore retains a small positive contribution
\(-\mu_a f_2^2\) in the mirror functional. A carrier derivative has
not been discarded as negligible before the algebraic cancellation.

\subsection{Payment for the fourth derivative}

Theorem~\ref{thm:local} uses an analytic error on a larger disk, so
there are explicit Cauchy payments
\[
 |f_j-q_j|\le\epsilon_j,\qquad
 \epsilon_j=j!R^{-j}C_{2R}a^{-3/4},\qquad 2\le j\le4,
\]
where \(q_j=\partial_w^jq_a(0)\).
All these jets are real. Expanding \eqref{eq:Jdef} gives
\begin{equation}\label{eq:jeterror}
\begin{split}
 |\cJ_a(f_a)-\cJ_a(q_a)|\le\Delta_{\rm jet}:={}&
 4|q_3|\epsilon_3+2\epsilon_3^2
 +3|q_2|\epsilon_4+3|q_4|\epsilon_2+3\epsilon_2\epsilon_4\\
 &+|\mu_a|(2|q_2|\epsilon_2+\epsilon_2^2).
\end{split}
\end{equation}
The moment estimate and the correction imply
\(|q_j|\le M_j+O(\chi)\). Therefore
\begin{equation}\label{eq:jeterrorsize}
 \Delta_{\rm jet}
 =O\left(a^{-3/4}\ee^{\delta L}/L^2+a^{-3/2}\right)=o(1).
\end{equation}
For the actual approximate field, the justified necessary statement is
\begin{equation}\label{eq:approxjetcondition}
 |q_a(0)|,\ |q_a'(0)|=O(a^{-3/4}),\qquad
 \cJ_a(q_a)\ge-\Delta_{\rm jet}.
\end{equation}
Exact nonnegativity of the approximate functional is not silently
substituted for this paid inequality.

\subsection{One point satisfies both signed tests}

\begin{theorem}[Joint mirror compatibility]\label{thm:joint}
Under the hypotheses of Theorem~\ref{thm:mult}, the same
multiplicative point can be chosen to satisfy
\begin{align}
 q_*(0)&=q_*'(0)=0,\qquad q_*''(0)<0,\qquad q_*^{(4)}(0)>0,\\
 \cD_a(q_*)&\le-\kappa_\beta\frac{\ee^{(1/2-c/8)L}}{L^8},
 \label{eq:jointD}\\
 \cJ_a(q_*)&\ge c_J\frac{\ee^{(1/2-c/8)L}}{L^8}>0,\qquad
 c_J=\frac{3h_2h_4}{25600}.
 \label{eq:jointJ}
\end{align}
In particular, Corollary~\ref{cor:growing}, including the opposite phase
at a prime \(q_0\in(Y,2Y]\), remains valid simultaneously.
\end{theorem}
\begin{proof}
Make the seven-observation choice in Lemma~\ref{lem:complement}, rather
than only the four-observation choice. The event repair and signed
estimate in Section~\ref{sec:multiplicative} are unchanged, with
constants still independent of \(\cQ\).

For each \(j=2,3,4\), the reserved phase repair has derivative cost
\[
 |(\Delta p)^{(j)}(0)|
 \le C W(|t|+|s|)L^{-j}
 =O_\beta(C_{\cQ}L^{1-j}).
\]
The chosen complementary derivatives are \(O_\beta(C_{\cQ})\);
the entire correction has derivatives \(O(\chi)\).
The budget \(\mathcal E_{\cQ}\to0\) implies
\(C_{\cQ}=o(W/L^4)\). Thus the errors are smaller than both leading
quantities in \eqref{eq:primejets}, and eventually
\begin{equation}\label{eq:finaljets}
 q_*''(0)\le-\tfrac12h_2W/L^2,\qquad
 q_*^{(4)}(0)\ge\tfrac12h_4W/L^4.
\end{equation}
All jets are real, and \(\mu_a<0\). In \eqref{eq:Jdef},
the \(2q_3^2\) and \(-\mu_a q_2^2\) terms are nonnegative.
The cross term therefore gives
\[
 \cJ_a(q_*)\ge-3q_*''(0)q_*^{(4)}(0)
        \ge\tfrac34h_2h_4W^2/L^6.
\]
Use \eqref{eq:Wpopulation} to obtain \eqref{eq:jointJ}.
Equation~\eqref{eq:jointD} holds for this very same point by the
previous signed-error calculation.
\end{proof}

The model now has an exact local double zero at \(w=0\), since its
second derivative is nonzero. It passes the exact mirror functional,
which is stronger than merely passing the approximate necessary
condition \eqref{eq:approxjetcondition}. Nevertheless this is not
a physical collision: the model is not identified with a globally
even canonical field, and no all-real global zero configuration is
constructed. The opposite calibrated-prime defect explicitly forbids
the actual same-height phase law.

The conclusion is limited to the one mirror inequality tested here.
It does not establish compatibility with every finite collection of
Laguerre inequalities, every higher event jet, or global all-realness.
Each stronger constraint would need its own derivation, approximation
cost, and simultaneous construction or exclusion argument.
